\documentclass[12pt,a4paper]{article}
\usepackage[utf8]{inputenc}
\usepackage[spanish]{babel}
\usepackage{geometry}
\usepackage{fancyhdr}
\usepackage{enumitem}
\usepackage{xcolor}
\usepackage{listings}
\usepackage{graphicx}
\usepackage{hyperref}
\usepackage{tabularx}
\usepackage{array}

\geometry{margin=2cm}
\pagestyle{fancy}
\fancyhf{}
\lhead{Curso de Programación en Python}
\rhead{Semana 3}
\cfoot{\thepage}

% Configuración para código Python
\lstset{
    language=Python,
    basicstyle=\ttfamily\small,
    keywordstyle=\color{blue},
    stringstyle=\color{red},
    commentstyle=\color{green!60!black},
    showstringspaces=false,
    breaklines=true,
    frame=single,
    numbers=left,
    numberstyle=\tiny\color{gray}
}

\title{\textbf{Tarea Semana 3: Validador de Datos de Usuario}\\
\large Excepciones y Manejo de Errores}
\author{Curso de Programación en Python}
\date{Tiempo de entrega: 1 semana}

\begin{document}

\maketitle

\section{Introducción}

¡Bienvenido a la Semana 3! Ya sabes programar con variables, tomar decisiones con condicionales, y repetir acciones con bucles. Ahora aprenderás una habilidad crucial: \textbf{manejar errores de forma elegante}.

En el mundo real, los programas no siempre reciben datos perfectos. Los usuarios cometen errores, ingresan texto donde debería haber números, dividen entre cero, o intentan abrir archivos que no existen. Un buen programa no colapsa ante estos errores, sino que los maneja apropiadamente y guía al usuario hacia el uso correcto.

Las \textbf{excepciones} son el mecanismo de Python para manejar estos errores. Con \texttt{try-except}, puedes anticipar problemas, capturarlos antes de que hagan crash tu programa, y responder de forma apropiada.

En esta tarea, crearás un sistema robusto de validación de datos que maneja todo tipo de errores con gracia y profesionalismo.

\section{Objetivos de Aprendizaje}

Al completar esta tarea, serás capaz de:

\begin{itemize}[leftmargin=*]
    \item Identificar tipos de errores comunes: \texttt{ValueError}, \texttt{ZeroDivisionError}, \texttt{TypeError}, \texttt{IndexError}
    \item Usar bloques \texttt{try-except} para capturar y manejar excepciones
    \item Implementar la cláusula \texttt{else} en manejo de excepciones
    \item Usar la cláusula \texttt{finally} para código que siempre debe ejecutarse
    \item Crear excepciones personalizadas con \texttt{raise}
    \item Manejar múltiples tipos de excepciones en un solo bloque
    \item Validar entradas de usuario de forma robusta
    \item Escribir código que nunca haga crash inesperadamente
    \item Dar mensajes de error útiles y comprensibles
\end{itemize}

\section{Enunciado del Proyecto}

\subsection{Descripción General}

Vas a crear un programa llamado \textbf{``Sistema de Registro de Estudiantes de la UCR''}. Este sistema permite registrar información de estudiantes costarricenses y validar TODOS los datos de forma exhaustiva usando manejo de excepciones.

El sistema debe capturar errores, dar mensajes claros, y nunca terminar inesperadamente. Cada entrada debe validarse múltiples veces hasta que sea correcta.

\subsection{Datos a Registrar}

Para cada estudiante se debe capturar:

\begin{enumerate}
    \item \textbf{Información Personal:}
    \begin{itemize}
        \item Nombre completo (no vacío, solo letras y espacios)
        \item Edad (número entero entre 17 y 99)
        \item Cédula (formato: 1-2345-6789, validación de formato)
        \item Correo electrónico (debe contener @ y un dominio)
    \end{itemize}
    
    \item \textbf{Información Académica:}
    \begin{itemize}
        \item Carrera (no vacía, de una lista predefinida)
        \item Año de ingreso (entre 2015 y 2025)
        \item Promedio actual (decimal entre 0.0 y 100.0)
        \item Cantidad de créditos aprobados (entero no negativo)
    \end{itemize}
    
    \item \textbf{Información Financiera:}
    \begin{itemize}
        \item Tipo de beca (ninguna, media, completa)
        \item Monto de matrícula (decimal positivo)
        \item Estado de pago (al día / pendiente)
    \end{itemize}
\end{enumerate}

\section{Requisitos Funcionales}

\subsection{Parte 1: Validaciones con Try-Except}

\subsubsection{Validación de Edad}

Debe manejar los siguientes errores:

\begin{lstlisting}[caption=Validación robusta de edad]
def validar_edad():
    while True:
        try:
            edad_str = input("Ingresa tu edad: ")
            
            # Convertir a entero (puede lanzar ValueError)
            edad = int(edad_str)
            
            # Validar rango con raise
            if edad < 17:
                raise ValueError("Debes tener al menos 17 años")
            if edad > 99:
                raise ValueError("Edad invalida (maximo 99)")
            
            # Si llega aqui, la edad es valida
            return edad
            
        except ValueError as e:
            # Capturar tanto errores de conversion como personalizados
            if "invalid literal" in str(e):
                print("Error: Debes ingresar un numero entero")
            else:
                print(f"Error: {e}")
        except Exception as e:
            print(f"Error inesperado: {e}")
\end{lstlisting}

\subsubsection{Validación de Nombre}

\begin{lstlisting}[caption=Validación de nombre con excepciones]
def validar_nombre():
    while True:
        try:
            nombre = input("Nombre completo: ").strip()
            
            # Validar que no este vacio
            if not nombre:
                raise ValueError("El nombre no puede estar vacio")
            
            # Validar que solo contenga letras y espacios
            if not all(c.isalpha() or c.isspace() for c in nombre):
                raise ValueError("El nombre solo puede contener letras")
            
            # Validar longitud minima
            if len(nombre) < 3:
                raise ValueError("El nombre es demasiado corto")
            
            return nombre.title()
            
        except ValueError as e:
            print(f"Error: {e}")
\end{lstlisting}

\subsection{Parte 2: Validaciones Complejas}

\subsubsection{Validación de Cédula Costarricense}

Formato: X-XXXX-XXXX (donde X son dígitos)

\begin{lstlisting}[caption=Validación de cedula]
def validar_cedula():
    while True:
        try:
            cedula = input("Cedula (formato 1-2345-6789): ").strip()
            
            # Validar formato basico
            if not cedula:
                raise ValueError("La cedula no puede estar vacia")
            
            # Separar por guiones
            partes = cedula.split('-')
            
            # Validar cantidad de partes
            if len(partes) != 3:
                raise ValueError("Formato incorrecto. Usa: X-XXXX-XXXX")
            
            # Validar longitud de cada parte
            if len(partes[0]) != 1:
                raise ValueError("Primera parte debe tener 1 digito")
            if len(partes[1]) != 4:
                raise ValueError("Segunda parte debe tener 4 digitos")
            if len(partes[2]) != 4:
                raise ValueError("Tercera parte debe tener 4 digitos")
            
            # Validar que sean numeros
            for parte in partes:
                if not parte.isdigit():
                    raise ValueError("La cedula solo debe contener digitos")
            
            return cedula
            
        except ValueError as e:
            print(f"Error: {e}")
        except IndexError:
            print("Error: Formato de cedula invalido")
\end{lstlisting}

\subsubsection{Validación de Correo}

\begin{lstlisting}[caption=Validación basica de correo electronico]
def validar_correo():
    while True:
        try:
            correo = input("Correo electronico: ").strip().lower()
            
            if not correo:
                raise ValueError("El correo no puede estar vacio")
            
            # Validar que contenga @
            if '@' not in correo:
                raise ValueError("El correo debe contener @")
            
            # Separar usuario y dominio
            partes = correo.split('@')
            
            if len(partes) != 2:
                raise ValueError("El correo debe tener exactamente un @")
            
            usuario, dominio = partes
            
            if not usuario:
                raise ValueError("El correo debe tener usuario antes del @")
            
            if not dominio:
                raise ValueError("El correo debe tener dominio despues del @")
            
            # Validar que el dominio tenga punto
            if '.' not in dominio:
                raise ValueError("El dominio debe contener un punto")
            
            return correo
            
        except ValueError as e:
            print(f"Error: {e}")
        except Exception as e:
            print(f"Error inesperado: {e}")
\end{lstlisting}

\subsection{Parte 3: Uso de Try-Except-Else-Finally}

\subsubsection{Validación de Promedio con Else}

\begin{lstlisting}[caption=Uso de else en try-except]
def validar_promedio():
    while True:
        try:
            promedio_str = input("Promedio (0.0 - 100.0): ")
            promedio = float(promedio_str)
            
            if promedio < 0 or promedio > 100:
                raise ValueError("El promedio debe estar entre 0 y 100")
                
        except ValueError as e:
            if "could not convert" in str(e):
                print("Error: Debes ingresar un numero decimal")
            else:
                print(f"Error: {e}")
        else:
            # Este bloque solo se ejecuta si NO hubo excepcion
            print(f"Promedio {promedio:.2f} registrado exitosamente")
            return promedio
        finally:
            # Este bloque SIEMPRE se ejecuta
            print("Validacion de promedio completada")
\end{lstlisting}

\subsection{Parte 4: Cálculos con Manejo de Errores}

\subsubsection{Cálculo de Monto a Pagar (con ZeroDivisionError)}

\begin{lstlisting}[caption=Manejo de division entre cero]
def calcular_monto_pagar(matricula, creditos, tipo_beca):
    try:
        # Calcular descuento segun tipo de beca
        if tipo_beca == "completa":
            descuento = 1.0  # 100%
        elif tipo_beca == "media":
            descuento = 0.5  # 50%
        else:
            descuento = 0.0  # 0%
        
        # Calcular monto con descuento
        monto_descuento = matricula * descuento
        monto_final = matricula - monto_descuento
        
        # Calcular costo por credito
        if creditos == 0:
            raise ZeroDivisionError("No se puede calcular con 0 creditos")
        
        costo_por_credito = monto_final / creditos
        
        return monto_final, costo_por_credito
        
    except ZeroDivisionError as e:
        print(f"Error de division: {e}")
        return matricula, 0.0
    except TypeError as e:
        print(f"Error de tipo: {e}")
        return 0.0, 0.0
    finally:
        print("Calculo de montos finalizado")
\end{lstlisting}

\subsection{Parte 5: Menú Principal Robusto}

El programa debe tener un menú que NUNCA falle:

\begin{enumerate}[leftmargin=*]
    \item Registrar nuevo estudiante
    \item Ver estudiantes registrados
    \item Buscar estudiante por cédula
    \item Calcular estadísticas
    \item Salir
\end{enumerate}

Todo el menú debe estar envuelto en try-except para capturar cualquier error.

\newpage

\subsection{Ejemplo de Ejecución}

\begin{lstlisting}[caption=Flujo completo con manejo de errores]
========================================
  SISTEMA DE REGISTRO DE ESTUDIANTES
========================================

--- MENU PRINCIPAL ---
1. Registrar nuevo estudiante
2. Ver estudiantes registrados
3. Buscar estudiante
4. Estadisticas
5. Salir

Elige una opcion: abc
Error: Debes ingresar un numero entero

Elige una opcion: 9
Error: Opcion invalida. Debe ser entre 1 y 5

Elige una opcion: 1

--- REGISTRO DE ESTUDIANTE ---

=== INFORMACION PERSONAL ===

Nombre completo: 123
Error: El nombre solo puede contener letras

Nombre completo: An
Error: El nombre es demasiado corto

Nombre completo: Ana Maria Gomez
Validacion exitosa!

Ingresa tu edad: veinte
Error: Debes ingresar un numero entero

Ingresa tu edad: 15
Error: Debes tener al menos 17 años

Ingresa tu edad: 20
Validacion exitosa!

Cedula (formato 1-2345-6789): 12345678
Error: Formato incorrecto. Usa: X-XXXX-XXXX

Cedula (formato 1-2345-6789): 1-23-4567
Error: Segunda parte debe tener 4 digitos

Cedula (formato 1-2345-6789): 1-2345-6789
Validacion exitosa!

Correo electronico: ana.gmail.com
Error: El correo debe contener @

Correo electronico: ana@
Error: El correo debe tener dominio despues del @

Correo electronico: ana@gmail.com
Validacion exitosa!

=== INFORMACION ACADEMICA ===

Carreras disponibles:
1. Ingenieria en Computacion
2. Administracion de Empresas
3. Ingenieria en Diseno Industrial
4. Arquitectura
5. Ingenieria en Biotecnologia

Elige tu carrera (1-5): 6
Error: Opcion invalida

Elige tu carrera (1-5): 1
Carrera seleccionada: Ingenieria en Computacion

Año de ingreso (2015-2025): 2010
Error: Año invalido. Debe ser entre 2015 y 2025

Año de ingreso (2015-2025): 2022
Validacion exitosa!

Promedio (0.0 - 100.0): cien
Error: Debes ingresar un numero decimal
Validacion de promedio completada

Promedio (0.0 - 100.0): 150
Error: El promedio debe estar entre 0 y 100
Validacion de promedio completada

Promedio (0.0 - 100.0): 85.5
Promedio 85.50 registrado exitosamente
Validacion de promedio completada

Creditos aprobados: -10
Error: Los creditos no pueden ser negativos

Creditos aprobados: 45
Validacion exitosa!

=== INFORMACION FINANCIERA ===

Tipo de beca:
1. Ninguna
2. Media beca (50%)
3. Beca completa (100%)

Elige tipo de beca: 2
Beca seleccionada: Media beca

Monto de matricula: mil
Error: Debes ingresar un numero

Monto de matricula: 150000
Validacion exitosa!

Calculo de montos finalizado
Monto final a pagar: ₡75,000.00
Costo por credito: ₡1,666.67

--- RESUMEN DEL REGISTRO ---
Nombre: Ana Maria Gomez
Edad: 20 años
Cedula: 1-2345-6789
Correo: ana@gmail.com
Carrera: Ingenieria en Computacion
Año de ingreso: 2022
Promedio: 85.50
Creditos: 45
Beca: Media beca
Matricula original: ₡150,000.00
Monto a pagar: ₡75,000.00

Estudiante registrado exitosamente!
\end{lstlisting}

\newpage

\section{Requisitos Técnicos Detallados}

\subsection{Tipos de Excepciones a Manejar}

Tu programa DEBE capturar y manejar:

\begin{enumerate}[leftmargin=*]
    \item \textbf{ValueError:}
    \begin{itemize}
        \item Conversión de string a int/float falla
        \item Valores fuera de rango válido
        \item Excepciones personalizadas con \texttt{raise ValueError()}
    \end{itemize}
    
    \item \textbf{ZeroDivisionError:}
    \begin{itemize}
        \item División entre cero en cálculos
        \item Debe capturarse y manejar apropiadamente
    \end{itemize}
    
    \item \textbf{IndexError:}
    \begin{itemize}
        \item Acceso a índices inválidos en listas
        \item Por ejemplo, al parsear formatos
    \end{itemize}
    
    \item \textbf{TypeError:}
    \begin{itemize}
        \item Operaciones con tipos incompatibles
        \item Intentar concatenar string con número, etc.
    \end{itemize}
    
    \item \textbf{Exception (genérica):}
    \begin{itemize}
        \item Como último recurso para errores inesperados
        \item Debe ser el último except en la cadena
    \end{itemize}
\end{enumerate}

\subsection{Estructura de Try-Except-Else-Finally}

\textbf{Uso obligatorio en al menos 3 validaciones:}

\begin{lstlisting}
try:
    # Codigo que puede fallar
    valor = int(input("Ingresa un numero: "))
    
except ValueError:
    # Manejar error especifico
    print("Error: Debes ingresar un numero")
    
except Exception as e:
    # Manejar cualquier otro error
    print(f"Error inesperado: {e}")
    
else:
    # Solo se ejecuta si NO hubo excepcion
    print("Valor ingresado correctamente")
    
finally:
    # SIEMPRE se ejecuta, haya o no haya excepcion
    print("Validacion finalizada")
\end{lstlisting}

\subsection{Uso de Raise para Excepciones Personalizadas}

Debes usar \texttt{raise} para validaciones lógicas:

\begin{lstlisting}
def validar_rango(valor, minimo, maximo):
    try:
        if valor < minimo:
            raise ValueError(f"Valor muy bajo (minimo: {minimo})")
        if valor > maximo:
            raise ValueError(f"Valor muy alto (maximo: {maximo})")
        return True
    except ValueError as e:
        print(f"Error: {e}")
        return False
\end{lstlisting}

\subsection{Múltiples Excepciones en un Bloque}

Puedes capturar múltiples excepciones:

\begin{lstlisting}
try:
    # Codigo riesgoso
    resultado = operacion_compleja()
except (ValueError, TypeError) as e:
    # Manejar multiples tipos de excepciones
    print(f"Error de valor o tipo: {e}")
except ZeroDivisionError:
    # Manejar division por cero especificamente
    print("Error: Division entre cero")
\end{lstlisting}

\section{Consejos y Recomendaciones}

\begin{itemize}[leftmargin=*]
    \item \textbf{No uses try-except para todo:} Solo usa donde realmente pueda haber un error.
    
    \item \textbf{Sé específico con las excepciones:} Captura el tipo exacto de error cuando sea posible.
    
    \item \textbf{Mensajes de error claros:} Explica qué salió mal y cómo corregirlo.
    
    \item \textbf{No escondas errores:} Un \texttt{except} vacío o que solo hace \texttt{pass} es mala práctica.
    
    \item \textbf{Valida antes de convertir:} Usa métodos como \texttt{.isdigit()} antes de \texttt{int()}.
    
    \item \textbf{Usa raise para validaciones:} Es más claro que retornar valores especiales.
    
    \item \textbf{Finally para limpieza:} Usa \texttt{finally} para cerrar archivos, conexiones, etc.
    
    \item \textbf{Prueba todos los errores:} Ingresa datos incorrectos intencionalmente para verificar el manejo.
    
    \item \textbf{Evita try-except muy amplios:} Agrupa solo el código que puede fallar.
\end{itemize}

\section{Entregables}

\begin{enumerate}[leftmargin=*]
    \item Archivo Python llamado \texttt{registro\_estudiantes.py}
    \item El código NO debe crashear bajo ninguna circunstancia
    \item Debe manejar TODAS las validaciones con try-except
    \item Debe incluir comentarios explicando el manejo de errores
    \item Debe demostrar uso de try-except-else-finally
    \item Debe usar \texttt{raise} para excepciones personalizadas
\end{enumerate}

\section{Defensa del Código (60\% de la nota final)}

\subsection{¿Qué es la defensa del código?}

La defensa del código es una reunión virtual donde presentarás y explicarás tu trabajo. Esta defensa es \textbf{obligatoria} y representa el \textbf{60\% de la nota final} de esta tarea.

\subsection{¿Cómo funciona?}

Durante la defensa deberás:

\begin{enumerate}[leftmargin=*]
    \item \textbf{Compartir tu pantalla} mostrando tu código en el editor.
    
    \item \textbf{Demostrar el manejo de errores:}
    \begin{itemize}
        \item Ingresar datos incorrectos intencionalmente
        \item Mostrar que el programa NUNCA crashea
        \item Demostrar cada tipo de excepción capturada
        \item Mostrar validaciones funcionando correctamente
    \end{itemize}
    
    \item \textbf{Explicar tus bloques try-except:}
    \begin{itemize}
        \item Por qué elegiste capturar cada tipo de excepción
        \item Cómo funciona cada bloque de validación
        \item Qué hace el código en \texttt{else} vs \texttt{finally}
        \item Cuándo y por qué usas \texttt{raise}
    \end{itemize}
    
    \item \textbf{Responder preguntas como:}
    \begin{itemize}
        \item ``¿Qué diferencia hay entre \texttt{ValueError} y \texttt{TypeError}?''
        \item ``¿Por qué usaste \texttt{raise} aquí?''
        \item ``¿Qué pasaría si no capturaras este error?''
        \item ``¿Cuál es la diferencia entre \texttt{else} y poner el código después del try?''
        \item ``¿Cuándo se ejecuta el bloque \texttt{finally}?''
        \item ``¿Qué hace \texttt{Exception as e}?''
    \end{itemize}
    
    \item \textbf{Demostrar comprensión de:}
    \begin{itemize}
        \item Diferentes tipos de excepciones
        \item Cuándo usar try-except vs validación con if
        \item El flujo de ejecución con try-except-else-finally
        \item Buenas prácticas en manejo de errores
    \end{itemize}
\end{enumerate}

\subsection{Preguntas Típicas de la Defensa}

Prepárate para responder:

\begin{itemize}[leftmargin=*]
    \item ¿Cuál es la diferencia entre \texttt{except ValueError} y \texttt{except Exception}?
    \item ¿Por qué el orden de los \texttt{except} importa?
    \item ¿Cuándo usarías \texttt{raise} en lugar de \texttt{print}?
    \item ¿Qué significa \texttt{as e} en \texttt{except ValueError as e}?
    \item ¿Cuál es la diferencia entre \texttt{else} y \texttt{finally}?
    \item Dame un ejemplo de cuándo usar \texttt{ZeroDivisionError}
    \item ¿Por qué no usar \texttt{except:} sin especificar el tipo?
    \item ¿Cómo probaste que tu programa maneja todos los errores?
\end{itemize}

\subsection{Consejos para una Buena Defensa}

\begin{itemize}[leftmargin=*]
    \item \textbf{Conoce cada excepción:} Entiende qué causa cada tipo de error
    \item \textbf{Practica romper tu programa:} Ingresa datos inválidos y verifica el comportamiento
    \item \textbf{Explica el flujo:} Describe qué pasa cuando hay y cuando no hay error
    \item \textbf{Ten ejemplos preparados:} De cada tipo de excepción que capturas
    \item \textbf{Entiende raise:} Sé claro sobre cuándo y por qué crear excepciones
\end{itemize}

\subsection{Distribución de la Nota}

\begin{itemize}
    \item \textbf{Código funcional (40\%):} El programa maneja errores sin crashear
    \item \textbf{Defensa del código (60\%):} Tu capacidad de explicar el manejo de excepciones
\end{itemize}

\textbf{Importante:} Si no asistes a la defensa o no puedes explicar tus try-except, la nota máxima será 40/100.

\section{Desafíos Opcionales (Puntos Extra)}

\begin{enumerate}[leftmargin=*]
    \item \textbf{Registro en archivo (+7 pts):}
    \begin{itemize}
        \item Guardar estudiantes en archivo de texto
        \item Manejar \texttt{FileNotFoundError}, \texttt{PermissionError}
        \item Usar \texttt{try-except-finally} para cerrar archivos apropiadamente
    \end{itemize}
    
    \item \textbf{Validación avanzada de cédula (+5 pts):}
    \begin{itemize}
        \item Verificar dígito verificador real de cédulas costarricenses
        \item Validar provincia según primer dígito
        \item Manejar excepciones personalizadas
    \end{itemize}
    
    \item \textbf{Sistema de log de errores (+6 pts):}
    \begin{itemize}
        \item Registrar en archivo todos los errores que ocurren
        \item Incluir timestamp, tipo de error, mensaje
        \item Usar módulo \texttt{datetime}
    \end{itemize}
    
    \item \textbf{Excepciones personalizadas (+5 pts):}
    \begin{itemize}
        \item Crear clases de excepciones propias: \texttt{CedulaInvalidaError}, \texttt{PromedioInvalidoError}
        \item Usar correctamente en el programa
        \item Explicar en la defensa cómo funcionan
    \end{itemize}
    
    \item \textbf{Reintentos limitados (+4 pts):}
    \begin{itemize}
        \item Permitir solo 3 intentos por campo
        \item Si falla 3 veces, preguntar si desea continuar o cancelar
        \item Usar excepciones para controlar el flujo
    \end{itemize}
\end{enumerate}

\textit{Nota: Los puntos extra solo se otorgan si la implementación es correcta Y puedes explicarla en la defensa.}

\newpage

\section{Rúbrica de Evaluación}

\subsection{Distribución General}

\begin{itemize}
    \item \textbf{Código Funcional:} 40\% (40 puntos)
    \item \textbf{Defensa del Código:} 60\% (60 puntos)
    \item \textbf{Total:} 100 puntos
    \item \textbf{Puntos Extra:} Hasta +27 puntos adicionales
\end{itemize}

\subsection{Parte 1: Código Funcional (40 puntos)}

\subsubsection{1.1 Implementación de Try-Except (20 puntos)}

\begin{itemize}[leftmargin=*]
    \item \textbf{Try-except básico (5 pts):}
    \begin{itemize}
        \item 5 pts: Usa try-except correctamente en múltiples validaciones
        \item 4 pts: Usa en varias validaciones con errores menores
        \item 2 pts: Uso mínimo o incorrecto
        \item 0 pts: No usa try-except
    \end{itemize}
    
    \item \textbf{Manejo de ValueError (4 pts):}
    \begin{itemize}
        \item 4 pts: Captura y maneja apropiadamente en conversiones y validaciones
        \item 3 pts: Maneja pero con mensajes poco claros
        \item 1 pt: Implementación parcial
        \item 0 pts: No maneja ValueError
    \end{itemize}
    
    \item \textbf{Manejo de múltiples tipos de excepciones (5 pts):}
    \begin{itemize}
        \item 5 pts: Maneja al menos 4 tipos diferentes de excepciones apropiadamente
        \item 4 pts: Maneja 3 tipos diferentes
        \item 2 pts: Maneja solo 1-2 tipos
        \item 0 pts: Solo maneja Exception genérica
    \end{itemize}
    
    \item \textbf{Uso de raise (3 pts):}
    \begin{itemize}
        \item 3 pts: Usa raise correctamente para excepciones personalizadas
        \item 2 pts: Usa raise pero no en todos los casos apropiados
        \item 1 pt: Uso mínimo o incorrecto
        \item 0 pts: No usa raise
    \end{itemize}
    
    \item \textbf{Try-except-else-finally (3 pts):}
    \begin{itemize}
        \item 3 pts: Usa la estructura completa en al menos 2 validaciones
        \item 2 pts: Usa en 1 validación correctamente
        \item 1 pt: Usa pero incorrectamente
        \item 0 pts: No usa else ni finally
    \end{itemize}
\end{itemize}

\subsubsection{1.2 Validaciones Implementadas (12 puntos)}

\begin{itemize}[leftmargin=*]
    \item \textbf{Validación de información personal (4 pts):}
    \begin{itemize}
        \item 4 pts: Valida nombre, edad, cédula y correo con try-except
        \item 3 pts: Valida 3 de 4 campos
        \item 2 pts: Valida 2 de 4 campos
        \item 0 pts: Validación mínima o sin excepciones
    \end{itemize}
    
    \item \textbf{Validación de información académica (4 pts):}
    \begin{itemize}
        \item 4 pts: Valida carrera, año, promedio y créditos apropiadamente
        \item 3 pts: Valida 3 de 4 campos
        \item 2 pts: Valida 2 de 4 campos
        \item 0 pts: Validación insuficiente
    \end{itemize}
    
    \item \textbf{Validación de información financiera (4 pts):}
    \begin{itemize}
        \item 4 pts: Valida tipo de beca y monto con manejo de errores
        \item 3 pts: Validación parcial
        \item 1 pt: Validación mínima
        \item 0 pts: No valida apropiadamente
    \end{itemize}
\end{itemize}

\subsubsection{1.3 Robustez del Programa (8 puntos)}

\begin{itemize}[leftmargin=*]
    \item \textbf{El programa NUNCA crashea (4 pts):}
    \begin{itemize}
        \item 4 pts: Imposible hacer crash el programa con entradas inválidas
        \item 3 pts: Maneja la mayoría de casos, falla en casos extremos
        \item 1 pt: Crashea con entradas inválidas comunes
        \item 0 pts: Crashea frecuentemente
    \end{itemize}
    
    \item \textbf{Mensajes de error claros (2 pts):}
    \begin{itemize}
        \item 2 pts: Todos los mensajes son claros y útiles
        \item 1 pt: Mensajes básicos pero comprensibles
        \item 0 pts: Mensajes confusos o genéricos
    \end{itemize}
    
    \item \textbf{Flujo de validación completo (2 pts):}
    \begin{itemize}
        \item 2 pts: Todas las entradas se validan en bucles hasta que sean correctas
        \item 1 pt: La mayoría de entradas se validan apropiadamente
        \item 0 pts: Validación incompleta o acepta datos inválidos
    \end{itemize}
\end{itemize}

\subsection{Parte 2: Defensa del Código (60 puntos)}

\subsubsection{2.1 Comprensión de Excepciones (30 puntos)}

\begin{itemize}[leftmargin=*]
    \item \textbf{Explicación de tipos de excepciones (10 pts):}
    \begin{itemize}
        \item 10 pts: Explica perfectamente ValueError, TypeError, ZeroDivisionError, etc.
        \item 7 pts: Explica correctamente la mayoría
        \item 4 pts: Explicación básica o con errores
        \item 0 pts: No entiende los tipos de excepciones
    \end{itemize}
    
    \item \textbf{Explicación de try-except-else-finally (10 pts):}
    \begin{itemize}
        \item 10 pts: Explica claramente cuándo se ejecuta cada bloque
        \item 7 pts: Explica correctamente con pequeñas imprecisiones
        \item 4 pts: Explicación superficial o confusa
        \item 0 pts: No entiende la estructura
    \end{itemize}
    
    \item \textbf{Explicación de raise (6 pts):}
    \begin{itemize}
        \item 6 pts: Explica perfectamente cuándo y por qué usar raise
        \item 4 pts: Explicación correcta pero incompleta
        \item 2 pts: Explicación vaga
        \item 0 pts: No entiende raise
    \end{itemize}
    
    \item \textbf{Buenas prácticas (4 pts):}
    \begin{itemize}
        \item 4 pts: Explica por qué no usar except: vacío, importancia de ser específico
        \item 3 pts: Conoce algunas buenas prácticas
        \item 1 pt: Conocimiento mínimo
        \item 0 pts: No conoce buenas prácticas
    \end{itemize}
\end{itemize}

\subsubsection{2.2 Demostración y Razonamiento (20 puntos)}

\begin{itemize}[leftmargin=*]
    \item \textbf{Demostración de manejo de errores (10 pts):}
    \begin{itemize}
        \item 10 pts: Demuestra múltiples tipos de errores siendo capturados correctamente
        \item 7 pts: Demuestra varios casos de error
        \item 4 pts: Demuestra casos básicos
        \item 0 pts: No puede demostrar el manejo de errores
    \end{itemize}
    
    \item \textbf{Respuesta a preguntas técnicas (10 pts):}
    \begin{itemize}
        \item 10 pts: Responde correcta y claramente todas las preguntas sobre excepciones
        \item 7 pts: Responde bien la mayoría
        \item 4 pts: Respuestas parciales o confusas
        \item 2 pts: Dificultad con preguntas básicas
        \item 0 pts: No puede responder sobre su código
    \end{itemize}
\end{itemize}

\subsubsection{2.3 Calidad del Código (10 puntos)}

\begin{itemize}[leftmargin=*]
    \item \textbf{Comentarios sobre manejo de errores (4 pts):}
    \begin{itemize}
        \item 4 pts: Comentarios claros explicando cada bloque try-except
        \item 3 pts: Comentarios presentes pero básicos
        \item 1 pt: Comentarios mínimos
        \item 0 pts: Sin comentarios
    \end{itemize}
    
    \item \textbf{Organización del código (3 pts):}
    \begin{itemize}
        \item 3 pts: Código perfectamente organizado, funciones bien estructuradas
        \item 2 pts: Código organizado con errores menores
        \item 1 pt: Poco organizado
        \item 0 pts: Código desorganizado
    \end{itemize}
    
    \item \textbf{Mensajes al usuario (3 pts):}
    \begin{itemize}
        \item 3 pts: Mensajes profesionales, claros y útiles en todos los errores
        \item 2 pts: Mensajes aceptables
        \item 1 pt: Mensajes básicos
        \item 0 pts: Mensajes confusos o ausentes
    \end{itemize}
\end{itemize}

\subsection{Puntos Extra (Hasta +27 puntos)}

\begin{itemize}[leftmargin=*]
    \item \textbf{Registro en archivo (+7 pts):} Guarda/lee estudiantes con manejo de FileNotFoundError
    \item \textbf{Validación avanzada de cédula (+5 pts):} Valida dígito verificador y provincia
    \item \textbf{Sistema de log de errores (+6 pts):} Registra todos los errores en archivo con timestamp
    \item \textbf{Excepciones personalizadas (+5 pts):} Crea y usa clases de excepciones propias
    \item \textbf{Reintentos limitados (+4 pts):} Máximo 3 intentos por campo con opción de cancelar
\end{itemize}

\textbf{Nota:} Los puntos extra solo se otorgan si la implementación es correcta Y el estudiante puede explicarla completamente durante la defensa.

\subsection{Penalizaciones}

\begin{itemize}[leftmargin=*]
    \item \textbf{-100\% de la defensa:} No asiste sin justificación (nota máxima = 40/100)
    \item \textbf{-50\% de la defensa:} Evidencia de no entender el manejo de excepciones
    \item \textbf{-10 pts:} El programa crashea con entradas inválidas comunes
    \item \textbf{-5 pts:} Entrega tardía de 1 día
    \item \textbf{-10 pts:} Entrega tardía de 2-3 días
    \item \textbf{-20 pts:} Entrega tardía de 4+ días
    \item \textbf{-3 pts:} Nombre de archivo incorrecto
    \item \textbf{-5 pts:} Usa \texttt{except:} sin especificar tipo (mala práctica)
    \item \textbf{-3 pts:} No usa raise en ninguna validación
\end{itemize}

\subsection{Nota Mínima Aprobatoria}

Para aprobar esta tarea necesitas obtener al menos \textbf{60 puntos de 100}.

\subsection{Escala de Calificación}

\begin{itemize}
    \item \textbf{90-100+ pts:} Excelente - Dominio completo de manejo de excepciones
    \item \textbf{80-89 pts:} Muy Bueno - Comprensión sólida de try-except
    \item \textbf{70-79 pts:} Bueno - Comprensión aceptable, necesita práctica
    \item \textbf{60-69 pts:} Suficiente - Comprensión básica, necesita refuerzo
    \item \textbf{0-59 pts:} Insuficiente - No aprobado
\end{itemize}

\section{Recursos de Apoyo}

\begin{itemize}[leftmargin=*]
    \item \textbf{Documentación oficial:} \url{https://docs.python.org/es/3/tutorial/errors.html}
    \item \textbf{Lista de excepciones:} \url{https://docs.python.org/es/3/library/exceptions.html}
    \item \textbf{Try-except:} Revisa las notas sobre captura de excepciones
    \item \textbf{Raise:} Practica crear excepciones personalizadas
    \item \textbf{Finally:} Entiende cuándo usar para limpieza de recursos
\end{itemize}

\section{Preguntas Frecuentes}

\textbf{P: ¿Cuál es la diferencia entre ValueError y TypeError?}\\
R: ValueError es cuando el tipo es correcto pero el valor no (ej: int("abc")). TypeError es cuando el tipo es incorrecto para la operación (ej: "5" + 5).

\textbf{P: ¿Cuándo debo usar raise?}\\
R: Cuando detectas un error lógico en tu programa que no es un error de Python sino de validación de negocio.

\textbf{P: ¿Cuál es la diferencia entre else y poner código después del try?}\\
R: El código en else solo se ejecuta si NO hubo excepción. El código después se ejecuta siempre (si no hay return/break).

\textbf{P: ¿Cuándo se ejecuta finally?}\\
R: SIEMPRE, haya o no haya excepción, incluso si hay return en try o except.

\textbf{P: ¿Por qué no usar except: sin especificar?}\\
R: Porque captura TODAS las excepciones, incluso las del sistema, y puede esconder errores importantes.

\textbf{P: ¿Puedo tener múltiples except para un try?}\\
R: Sí! De hecho es recomendado para manejar diferentes errores de forma diferente.

\textbf{P: ¿Es obligatorio validar TODO con try-except?}\\
R: No. Usa if para validaciones simples. Usa try-except para operaciones que pueden fallar (conversiones, división, etc.).

\vspace{1cm}

\begin{center}
\textbf{¡Éxito en tu sistema de validación!}\\
\textit{El manejo de excepciones es lo que separa un programa amateur\\
de uno profesional y robusto.}
\end{center}

\end{document}
